\begin{figure}[h!]
    \centering
    \includegraphics[width=0.7\linewidth]{../results/graphics_general/policy/mar_hydrogen_strategy_semistacked.pdf}
    \caption{Hydrogen Strategy of Morocco \cite{MarHyStrat2021}. The export volume rises up to 115~TWh in 2050, the domestic demand up to 40 TWh in 2050. The largest share of domestic demand is in the industry sector. The exports are dominated by hydrogen, followed by synthetic fuels.}
    \label{fig:mar_hydrogen_strategy}
\end{figure}


\begin{figure}[h!]
    \centering
    \includegraphics[width=0.7\linewidth]{../results/graphics_general/morocco_em.pdf}
    \caption{Morocco's historical \co and greenhouse gas emissions (GHG). While emissions are rising, emission targets are in the range between 75 Mt\coe\ (conditional) and 115 Mt\coe\ (unconditional)\cite{CAT2021}. GHG = greenhouse gases.}
    \label{fig:morocco_em}
\end{figure}

\begin{figure*}[h!]
    \centering
    \includegraphics[trim={0cm 8cm 0cm 8cm}, clip, width=\linewidth]{figuremerge/figure_three.pdf}
    \caption{Electricity supply and demand at fixed export levels and increasing domestic \co mitigation and vice versa with hourly temporal hydrogen regulation.
    Panel \textbf{a} shows the electricity supply and demand at fixed export (1 TWh) and 0--100\% domestic \co mitigation.
    Increasing domestic \co mitigation first phases out carbon-intensive coal generation in favor of combined-cycle gas turbines (CCGT), at medium to high domestic \co mitigation the electricity system is fully renewable supported by flexibility through Vehicle-to-Grid (V2G) and sector coupling. Increasing electricity demands include Battery Electric Vehicles (BEV) and hydrogen ({H${}_2${}}) generation for other sectors.
    Panel \textbf{b} shows the electricity supply and demand at fixed  domestic \co mitigation (0\%) and 1--120 TWh export.    
    At increasing hydrogen exports the additional electricity required for hydrogen electrolysis is covered by onshore wind and solar photovoltaics (PV), as imposed by the temporal hydrogen regulation. 
    See Supplementary Fig. 4 for electricity supply and demand if no hydrogen regulation is in place and Supplementary Fig. 9 for hydrogen exports up to 200 TWh. DAC = direct air capture. %\ref{fig:balances-ac-noreg}, \ref{fig:balances-ac-200}
    }
    \label{fig:balances-ac}
\end{figure*}

\begin{figure*}[h!]
    \centering

    \includegraphics[trim={0cm 7cm 0cm 7cm}, clip, width=\linewidth]{figuremerge/figure_four.pdf}
    \caption{
    The effects of hydrogen exports on domestic electricity consumers and domestic \co mitigation on hydrogen exporters.
    Panel \textbf{a} shows the effect of hydrogen ({H${}_2${}}) exports on domestic electricity cost with hourly hydrogen regulation. Therefore, the effect of hydrogen exports is isolated by normalizing the costs to 1 TWh hydrogen export in each column. This approach allows for a vertical interpretation in Panel \textbf{a} only, displaying the effects of hydrogen exports (vertical) at a certain domestic \co mitigation level.
    Panel \textbf{b} shows the effect of domestic \co mitigation on hydrogen export cost with hourly hydrogen regulation. Therefore, the effect of domestic \co mitigation is isolated by normalizing the costs to 0\% \co mitigation in each row. This approach allows for a horizontal interpretation in Figure Panel \textbf{b} only, displaying the effects of domestic \co mitigation (horizontal) at a certain hydrogen export quantity.
    Domestic electricity consumers profit from increasing hydrogen exports, especially at low domestic \co mitigation and high exports. Hydrogen exporters profit from domestic \co mitigation at medium mitigation efforts.
    See Supplementary Fig. 17 for absolute market-based costs for domestic electricity consumers and hydrogen exporters, Supplementary Fig. 6 for the interplay of hydrogen exports and domestic \co mitigation if no hydrogen regulation is in place and Supplementary Fig. 10 for hydrogen exports up to 200 TWh.
    } %\ref{fig:expenses_abs_120}, \ref{fig:expenses_default_120-noreg}, \ref{fig:expenses_default_200}
    \label{fig:expenses_default_120}
\end{figure*}


\begin{figure}[h!]
    \centering
    \includegraphics[trim={0cm 5cm 0cm 4cm}, clip, width=0.7\linewidth]{figuremerge/figure_five.pdf}
    \caption{
    The effect of temporal hydrogen regulation on electricity prices.
    Figure \ref{fig:pdc-120-0} shows the price duration curve at 120 TWh export and 0\% domestic \co mitigation. Stricter temporal hydrogen regulation pushes the price duration curve towards the left, as additional renewable electricity capacities phase out fossil generation with higher short-term marginal costs than renewables. Negative prices below -50 € per megawatt-hour (\eurMWh) are cut off.}
    \label{fig:pdc-120-0}
\end{figure}



\begin{figure*}[h!]
    \centering
    \includegraphics[trim={0cm 9cm 0cm 8cm}, clip, width=\linewidth]{figuremerge/figure_six.pdf}
    \caption{
    The effect of temporal hydrogen regulation on the electricity system and consumer costs.
    Panel \textbf{a} shows the electricity dispatch and demand and Panel \textbf{b} shows the cost for consumers for various (hydrogen) temporal matching regimes in the 120 TWh export and 0\% \co mitigation scenario. 
    Stricter temporal matching decreases carbon-intensive electricity generation (coal \& gas) for hydrogen generation and even domestic electricity consumers (Panel \textbf{a}). Cost for export hydrogen generation increase to fulfill the temporal matching constraint, whereas domestic electricity consumers profit from stricter hydrogen regulation (Panel \textbf{b}).
    CHP = combined heat and power.
    }
    \label{fig:expenses_rule}
\end{figure*}


\begin{figure}[h!]
    \centering
    \includegraphics[trim={0cm 7cm 0cm 7cm}, clip, width=\linewidth]{figuremerge/figure_seven.pdf} 
    \caption{Relative change of electricity and hydrogen consumer cost depending on the temporal hydrogen regulation.     
    Violin plots indicate median (white dot), interquartile range (thick black bar), and 1.5x interquartile range (thin black bar), using kernel density estimation to show the distribution shape of the mitigation-export scenarios, grouped into three stylized transition strategies (1–3) reflecting different sequences of export and mitigation ambition.
    Domestic electricity consumers (Panel \textbf{a}) profit across all export and mitigation scenarios but most in the group of slow export and quick \co mitigation scenarios. Hydrogen exporters (Panel \textbf{b}) experience higher cost with stricter temporal hydrogen regulation. The temporal hydrogen ({H${}_2${}}) regulation regulates the welfare distribution between both groups. See Supplementary Fig. 11 for the sensitivity of the export quantities and Supplementary Fig. 3 for scenario grouping.} %\ref{fig:expenses-relative-200},  \ref{fig:scenario-sel}
    \label{fig:expenses-relative-120}
\end{figure}


\begin{figure*}[h!]
    \centering
    \includegraphics[width=0.7\linewidth]{../results/graphics_general/brownfield_capacities.pdf}
    \caption{Current capacities of electricity generation and distribution in Morocco. Morocco's electricity generation portfolio is currently dominated by fossil generation (coal and gas), includes some hydropower plants and increasing but still minor capacities of onshore wind and solar PV. The data is obtained from Parzen et al.\cite{Parzen2023} and the visualization is based on H\"orsch et al.\cite{Horsch2018}. Boundaries depicted are based on the Global Administrative Areas\cite{gadm2022} and are intended for illustrative purposes only, not implying territorial claims.
    }
    \label{fig:MAR_brownfield}
\end{figure*}
